\section{Результаты моделирования}

\subsection{Параметры модели}

Численное моделирование проведено для радиального подшипника скольжения геометрии~\varGeomLetter{} ($R = \varRmm$~мм, $c = \varCmkm$~мкм, $L = \varLmm$~мм) с поверхностным микрорельефом \varDimpleTypeGenitive{} типа. Основные параметры расчёта сведены в таблицу~\ref{tab:input_data}.

\begin{table}[H]
\small
\caption{Исходные данные для расчёта}
\label{tab:input_data}
\begin{tabularx}{\textwidth}{|X|c|c|}
\hline
\centering\textbf{Параметр} & \textbf{Значение} & \textbf{Размерность} \\
\hline
Радиус вала $R$ & \varRmm & мм \\
\hline
Радиальный зазор $c$ & \varCmkm & мкм \\
\hline
Длина подшипника $L$ & \varLmm & мм \\
\hline
Относительная длина $L/D$ & \varLD & -- \\
\hline
Полуось по оси $Z$, $a$ & \varDimpleAmm & мм \\
\hline
Полуось по оси $\varphi$, $b$ & \varDimpleBmm & мм \\
\hline
Глубина углубления $h_p$ & \varDimpleHpmkm & мкм \\
\hline
Относительная глубина $h_p/c$ & \varDimpleHpRatio & -- \\
\hline
Число углублений по окружности $N_\varphi$ & \varNphi & шт. \\
\hline
Число рядов по длине $N_Z$ & \varNZ & шт. \\
\hline
Область текстурирования & \varPhiRange & -- \\
\hline
Частота вращения $n$ & \varN & об/мин \\
\hline
Вязкость смазки $\eta$ & \varEta & Па$\cdot$с \\
\hline
Диапазон эксцентриситета $\varepsilon$ & 0{,}05--0{,}80 & -- \\
\hline
Размер сетки & $500 \times 500$ & -- \\
\hline
Параметр релаксации $\omega$ & 1{,}5 & -- \\
\hline
Критерий сходимости $\delta$ & $< 10^{-5}$ & -- \\
\hline
Макс. число итераций & 50000 & -- \\
\hline
\end{tabularx}
\end{table}

Для заданной частоты вращения линейная скорость вала определяется как $U = \pi D n / 60 = \pi \cdot 0{,}100 \cdot 2980 / 60 \approx 15{,}60$~м/с. Моделирование проведено при 15 значениях эксцентриситета в рассматриваемом диапазоне.

\subsection{Распределение зазора и давления}

На рисунке~\ref{fig:pressure_2d} показано распределение безразмерного давления $P(\varphi)$ в среднем сечении ($Z = 0$) при $\varepsilon = 0{,}6$ для гладкого подшипника и подшипника со сферическими углублениями.

\begin{figure}[H]
\centering
\includegraphics[width=\textwidth]{figures/pressure_2d.png}
\caption{Распределение давления $P(\varphi)$ при $Z=0$ для гладкого и текстурированного подшипников ($\varepsilon = 0{,}6$)}
\label{fig:pressure_2d}
\end{figure}

В случае гладкого подшипника давление плавно нарастает в сужающейся части зазора и достигает единственного максимума вблизи $\varphi \approx \pi$, после чего обращается в нуль вследствие кавитации. Зона положительного давления охватывает приблизительно половину окружности подшипника. Введение сферического микрорельефа приводит к появлению локальных осцилляций давления в области текстурирования: на входной кромке каждого углубления давление возрастает, на выходной~--- снижается. Тем не менее интегральный уровень давления для текстурированного подшипника оказывается выше, обеспечивая дополнительную несущую способность.

На рисунке~\ref{fig:clearance_2d} представлен профиль зазора $H(\varphi)$ в среднем сечении при $\varepsilon = 0{,}6$.

\begin{figure}[H]
\centering
\includegraphics[width=\textwidth]{figures/clearance_2d.png}
\caption{Распределение зазора $H(\varphi)$ при $Z=0$ для подшипника с \varDimpleTypeGenitive{} углублениями ($\varepsilon = 0{,}6$)}
\label{fig:clearance_2d}
\end{figure}

Зазор изменяется по закону $H = 1 + \varepsilon\cos\varphi$, на который накладываются локальные увеличения, обусловленные сферическими углублениями. Следует отметить, что для сферической шапки профиль углубления обладает симметричной кривизной относительно центра, что обеспечивает плавное сопряжение с основной поверхностью. Максимальная глубина углубления в безразмерных единицах составляет $\Delta H_{\max} = h_p/c = 0{,}2$.

Трёхмерные поля зазора и давления при $\varepsilon = 0{,}6$ приведены на рисунках~\ref{fig:clearance_3d} и~\ref{fig:pressure_3d}.

\begin{figure}[H]
\centering
\includegraphics[width=\textwidth]{figures/clearance_3d.png}
\caption{Трёхмерное распределение зазора $H(\varphi, Z)$ ($\varepsilon = 0{,}6$)}
\label{fig:clearance_3d}
\end{figure}

Регулярная структура микрорельефа хорошо различима на рисунке~\ref{fig:clearance_3d}: 8 углублений размещены по окружности в 11 рядах вдоль длины подшипника в области $90^\circ - 270^\circ$.

\begin{figure}[H]
\centering
\includegraphics[width=\textwidth]{figures/pressure_3d.png}
\caption{Трёхмерное распределение давления $P(\varphi, Z)$ ($\varepsilon = 0{,}6$)}
\label{fig:pressure_3d}
\end{figure}

Трёхмерное распределение давления (рисунок~\ref{fig:pressure_3d}) демонстрирует, что помимо основного максимума, определяемого макрогеометрией клинового зазора, на каждом углублении формируются локальные возмущения. При набегании потока на углубление зазор расширяется, а на выходной кромке~--- сужается, что порождает микрогидродинамический эффект. Благодаря ограничению Гюмбеля ($P \geq 0$) суммарный вклад каждого углубления в несущую способность остаётся положительным.

Распределение давления неоднородно по длине подшипника: максимум наблюдается в среднем сечении ($Z = 0$), а у торцов ($Z = \pm 1$) давление спадает до нуля из-за торцевых утечек. Эффект текстурирования проявляется в наибольшей степени вблизи среднего сечения, где фоновое давление велико, и ослабевает у периферийных рядов углублений.

\subsection{Сравнительный анализ интегральных показателей}

Результаты расчёта интегральных характеристик для гладкого подшипника и подшипника со сферическими углублениями при выбранных значениях эксцентриситета приведены в таблицах~\ref{tab:smooth_results} и~\ref{tab:textured_results}.

\begin{table}[H]
\small
\centering
\caption{Характеристики гладкого подшипника}
\label{tab:smooth_results}
\begin{tabular}{|c|c|c|c|}
\hline
$\boldsymbol{\varepsilon}$ & $\boldsymbol{F}$\textbf{, Н} & $\boldsymbol{\mu}$ & $\boldsymbol{Q}$\textbf{, л/с} \\
\hline
0{,}05 & 551{,}5 & 0{,}2245 & 0{,}6852 \\
\hline
0{,}16 & 1775{,}8 & 0{,}0707 & 0{,}6843 \\
\hline
0{,}26 & 3141{,}8 & 0{,}0412 & 0{,}6825 \\
\hline
0{,}37 & 4805{,}2 & 0{,}0282 & 0{,}6801 \\
\hline
0{,}48 & 6984{,}6 & 0{,}0207 & 0{,}6772 \\
\hline
0{,}59 & 10124{,}4 & 0{,}0157 & 0{,}6744 \\
\hline
0{,}69 & 15347{,}9 & 0{,}0118 & 0{,}6717 \\
\hline
0{,}80 & 25964{,}3 & 0{,}0085 & 0{,}6698 \\
\hline
\end{tabular}
\end{table}

\begin{table}[H]
\small
\centering
\caption{Характеристики подшипника с \varDimpleTypeGenitive{} углублениями}
\label{tab:textured_results}
\begin{tabular}{|c|c|c|c|}
\hline
$\boldsymbol{\varepsilon}$ & $\boldsymbol{F}$\textbf{, Н} & $\boldsymbol{\mu}$ & $\boldsymbol{Q}$\textbf{, л/с} \\
\hline
0{,}05 & 615{,}5 & 0{,}2004 & 0{,}6884 \\
\hline
0{,}16 & 1862{,}6 & 0{,}0671 & 0{,}6874 \\
\hline
0{,}26 & 3256{,}7 & 0{,}0395 & 0{,}6857 \\
\hline
0{,}37 & 4967{,}4 & 0{,}0271 & 0{,}6832 \\
\hline
0{,}48 & 7238{,}4 & 0{,}0199 & 0{,}6804 \\
\hline
0{,}59 & 10635{,}5 & 0{,}0148 & 0{,}6774 \\
\hline
0{,}69 & 16458{,}1 & 0{,}0109 & 0{,}6747 \\
\hline
0{,}80 & 29891{,}4 & 0{,}0073 & 0{,}6725 \\
\hline
\end{tabular}
\end{table}

С увеличением эксцентриситета нагрузочная способность $F$ растёт нелинейно для обеих конфигураций подшипника. Текстурированный подшипник демонстрирует более высокие значения $F$ во всём исследуемом диапазоне $\varepsilon$. При больших эксцентриситетах ($\varepsilon = 0{,}80$) эффект текстурирования особенно заметен: нагрузочная способность подшипника со сферическими углублениями (29891{,}4~Н) превосходит значение для гладкого (25964{,}3~Н) на 15{,}1\%.

Коэффициент трения $\mu$ убывает с ростом эксцентриситета; для текстурированного подшипника его значения ниже при всех $\varepsilon$. Объёмный расход смазки $Q$ слабо зависит от $\varepsilon$ и для текстурированного подшипника несколько выше вследствие увеличения объёма смазочного слоя в зонах углублений.

Следует отметить, что по сравнению с вариантом~5 (геометрия~A: $R = 35$~мм, $c = 50$~мкм, $L = 56$~мм, тот же тип микрорельефа~--- сферическая шапка) подшипник геометрии~B обладает существенно большей абсолютной нагрузочной способностью, что обусловлено увеличенными размерами ($D = 100$~мм) и длиной. Однако относительный прирост от текстурирования при $\varepsilon = 0{,}6$ для геометрии~B составляет $+4{,}85$\%, тогда как для геометрии~A данный показатель заметно выше. Это объясняется тем, что при бо́льшем радиальном зазоре ($c = 70$~мкм) относительное влияние углублений глубиной $h_p = 12$~мкм на поле давления оказывается менее выраженным.

Для количественного сопоставления характеристик гладкого и текстурированного подшипников в таблице~\ref{tab:comparison} приведены данные при $\varepsilon = 0{,}6$.

\begin{table}[H]
\small
\centering
\caption{Сравнение характеристик при $\varepsilon = 0{,}6$}
\label{tab:comparison}
\begin{tabular}{|l|c|c|c|}
\hline
\textbf{Характеристика} & \textbf{Гладкий} & \textbf{С углублениями} & $\boldsymbol{\Delta}$\textbf{, \%} \\
\hline
Нагрузка $F$, Н & 10685{,}6 & 11203{,}4 & $+4{,}85$ \\
\hline
Коэфф. трения $\mu$ & 0{,}01507 & 0{,}01426 & $-5{,}34$ \\
\hline
Расход $Q$, л/с & 0{,}6739 & 0{,}6770 & $+0{,}45$ \\
\hline
\end{tabular}
\end{table}

Применение сферических углублений повышает несущую способность подшипника на $4{,}85$\% благодаря микрогидродинамическому эффекту, формируемому на кромках углублений. Коэффициент трения уменьшается на $5{,}34$\% за счёт роста средней толщины смазочного слоя в зонах текстурирования. Расход смазки увеличивается всего на $0{,}45$\%, что не предъявляет дополнительных требований к системе маслоснабжения и одновременно содействует теплоотводу.

На рисунках~\ref{fig:load_vs_eps}--\ref{fig:flow_vs_eps} представлены графики зависимостей интегральных характеристик от эксцентриситета.

\begin{figure}[H]
\centering
\includegraphics[width=\textwidth]{figures/load_vs_epsilon.png}
\caption{Нагрузочная способность $F(\varepsilon)$ для гладкого и текстурированного подшипников}
\label{fig:load_vs_eps}
\end{figure}

Зависимость $F(\varepsilon)$ для текстурированного подшипника располагается выше кривой гладкого при всех значениях эксцентриситета. По мере роста $\varepsilon$ абсолютная разница нагрузочных способностей увеличивается: при $\varepsilon = 0{,}05$ она составляет около 64~Н, тогда как при $\varepsilon = 0{,}80$ достигает почти 3930~Н. Такое поведение обусловлено степенной зависимостью давления от толщины зазора ($P \sim h^{-3}$): при уменьшении зазора эффект микрорельефа нарастает.

\begin{figure}[H]
\centering
\includegraphics[width=\textwidth]{figures/friction_vs_epsilon.png}
\caption{Коэффициент трения $\mu(\varepsilon)$ для гладкого и текстурированного подшипников}
\label{fig:friction_vs_eps}
\end{figure}

Коэффициент трения текстурированного подшипника остаётся ниже гладкого во всём диапазоне $\varepsilon$. Обе зависимости монотонно убывают, стремясь к нулю при $\varepsilon \to 1$. Максимальная абсолютная разница наблюдается при малых эксцентриситетах: при $\varepsilon = 0{,}05$ она равна $\Delta\mu = 0{,}024$, а при $\varepsilon = 0{,}80$~--- $\Delta\mu = 0{,}001$. Относительное снижение трения при этом сохраняется на сопоставимом уровне во всём диапазоне.

\begin{figure}[H]
\centering
\includegraphics[width=\textwidth]{figures/flow_vs_epsilon.png}
\caption{Расход смазки $Q(\varepsilon)$ для гладкого и текстурированного подшипников}
\label{fig:flow_vs_eps}
\end{figure}

Расход смазки текстурированного подшипника незначительно превышает расход гладкого аналога. Отклонение составляет менее 0{,}5\% и остаётся практически постоянным при изменении эксцентриситета.

Таким образом, нанесение сферических углублений обеспечивает одновременное повышение нагрузочной способности и снижение коэффициента трения при минимальном увеличении расхода смазки. Эффект текстурирования усиливается при росте эксцентриситета, что представляет практический интерес для тяжелонагруженных узлов.

\subsection*{Выводы по главе 3}
\phantomsection
\addcontentsline{toc}{subsection}{Выводы по главе 3}

Проведён расчёт характеристик радиального подшипника скольжения геометрии~B со сферическим микрорельефом (тип~5) при 15 значениях эксцентриситета. Построены поля давления и зазора, получены зависимости нагрузочной способности, коэффициента трения и расхода смазки от $\varepsilon$. Сравнение при $\varepsilon = 0{,}6$ выявило прирост нагрузочной способности на $4{,}85$\%, снижение коэффициента трения на $5{,}34$\% и увеличение расхода на $0{,}45$\%, что свидетельствует об эффективности сферического текстурирования для подшипников геометрии~B.
